%! TeX root = ../main.tex

\chapter{Estado del arte y marco de referencia}
En la literatura se busca detectar planitud a un nivel milimétrico utilizando técnicas avanzadas de óptica y matemáticas, sin embargo, la mayoría de los estudios realizados son con las láminas en movimiento, justo antes o durante el cortado.
Esto introduce una serie de retos, pues el mismo movimiento induce vibraciones que pueden pasar como defectos en la planitud, pero como vemos en estudios realizados por \textcite{Alonso2020, Alonso2021}, muestran que con una técnica llamada spline cúbica de Hermite, podemos eliminar este ruido de una forma muy elegante.
A pesar de que la mayoría de los artículos indagan en otras técnicas para la detección de planitud, tales como proyección de líneas láser \parencite{Molleda2011, Alonso2020, Alonso2021, Tsai2021, Zhu2026}, aquellos de especial interés serán los que integran técnicas de \textit{Time-of-Flight} \parencite{Lange2001, Du2026, Yoshida2026}.

\section{Marco teórico y fundamentos técnicos}
La tecnología \textit{Time-of-Flight} opera bajo el principio de una medición directa del tiempo que pasa entre la emisión de una señal luminosa y su recepción tras reflejarse en la superficie del objeto, tal y como se describe en \textcite{Lange2001}.

\begin{figure}[ht]
	\centering
	\documentclass[tikz, border=5pt]{standalone}
\usepackage{xcolor}
\usetikzlibrary{shapes,arrows.meta,calc}

\begin{document}
% Macro para el bloque 3D (toma el color base, p. ej. "gray")
\newcommand{\drawblock}[5]{
	\begin{scope}[shift={(#1)}, rotate=#2]
		\filldraw[fill=#3!30, draw=black] (0,0) rectangle (#4,#5);
		\filldraw[fill=#3!50, draw=black] (#4,0) -- ++(30:0.4) -- ++(0,#5) -- (#4,#5) -- cycle;
		\filldraw[fill=#3!70, draw=black] (0,#5) -- ++(30:0.4) -- ++(#4,0) -- (#4,#5) -- cycle;
	\end{scope}
}
\begin{tikzpicture}[>=Latex, font=\sffamily\fontsize{8pt}{10pt}\selectfont]

	% Coordenadas base
	\coordinate (O) at (0,0);
	\coordinate (laser_start) at (-3,0.3);
	\coordinate (splitter_center) at (-1.5,0.3);
	\coordinate (target_front) at (4.5,0.3);
	\coordinate (receiver_front) at (-1.5,-1.2);

	% Laser Box (3D) -> Pasas "gray" simple, sin signos de admiración
	\coordinate (laser_corner) at ($(laser_start)+(-2.1,-0.3)$);
	\drawblock{laser_corner}{0}{gray}{2.1}{0.6}
	\node[anchor=center, text width=1.5cm, align=center] at ($(laser_start)+(-1.05,0)$) {\bf Láser};

	% Tilted Beam Splitter
	\draw[rotate around={-60:(splitter_center)}, line width=2pt, draw=black!80]
	($(splitter_center)+(0,-0.6)$) -- ($(splitter_center)+(0,0.6)$);

	% Target Block (3D) -> Pasas "gray" simple
	\coordinate (target_corner) at ($(target_front)+(0,-1.2)$);
	\drawblock{target_corner}{0}{gray}{0.8}{2.4}
	\node[anchor=center, align=center] at ($(target_front)+(0.4,-1.8)$) {Superficie};
	\fill[black] (target_front) circle (0.15);

	% Receiver Block (3D) -> Pasas "gray" simple
	\coordinate (receiver_corner) at ($(receiver_front)+(-1.2,-0.8)$);
	\drawblock{receiver_corner}{0}{gray}{2.4}{0.8}
	\node[anchor=center, text width=2cm, align=center] at ($(receiver_front)+(0,-0.4)$) {\bf Receptor};

	% --- Trayectorias de haz ---
	\draw[red, line width=1.2pt, ->] (laser_start) -- (splitter_center);
	\draw[red, line width=1.2pt] (splitter_center) -- (target_front);

	\coordinate (return_start) at ($(target_front)+(0,-0.08)$);
	\coordinate (return_mid) at ($(splitter_center)+(0,-0.08)$);
	\draw[red, line width=1.2pt] (return_start) -- (return_mid);
	\draw[red, line width=1.2pt, ->] (return_start) -- ($(return_mid)!0.5!(return_start)$);

	\draw[red, line width=1.2pt, ->] (return_mid) -- (receiver_front);

	% --- Acotación (distancia R) ---
	\draw[thick] ($(splitter_center)+(0,0.6)$) -- ++(0,1.2);
	\draw[thick] ($(target_front)+(0,1.2)$) -- ++(0,0.6);
	\draw[<->, thick] ($(splitter_center)+(0,1.5)$) -- node[above, font=\bf] {R} ($(target_front)+(0,1.5)$);

\end{tikzpicture}
\end{document}

	\caption{Principio de medición por Tiempo de Vuelo (ToF).}
	\label{fig:tof_scheme}
\end{figure}
Como se ilustra en la figura \ref{fig:tof_scheme}, el sistema emplea una arquitectura óptica coaxial donde una fuente láser modulada emite un haz luminoso que atraviesa un divisor de haz (\textit{beam splitter}).
La señal llega a la superficie de la hoja de acero a una distancia radial \(R\).
La señal se refleja siguiendo la misma dirección y se desvía ortogonalmente por el divisor óptico hacia el sensor receptor.
Para cuantificar este recorrido, la luz emitida se modula con una señal periódica a una radiofrecuencia fija \(f_{\text{mod}}\) (normalmente de \qtyrange{20}{100}{\mega\hertz}).
El tiempo de vuelo \(\tau\) que tardan los fotones en completar el viaje provoca un retraso o desfase angular \(\varphi\) respecto a la señal de referencia.

Para extraer este valor, el receptor realiza cuatro capturas integradas de carga (\(A_0, A_1, A_2, A_3\)) con desfases programados de \ang{90} (\ang{0}, \ang{90}, \ang{180} y \ang{270}).
Con estas lecturas por pixel, el desfase \(\varphi(x, y)\) y la amplitud modulada \(B(x, y)\) se calculan directamente:

\begin{equation}
	\varphi(x, y) = \arctan\parens*{\frac{A_3(x, y) - A_1(x, y)}{A_0(x, y) - A_2(x, y)}}
\end{equation}

\begin{equation}
	B(x, y) = \frac{1}{2} \sqrt{\parens{A_3(x, y) - A_1(x, y)}^2 + \parens{A_0(x, y) - A_2(x, y)}^2}
\end{equation}

El valor de \(B(x, y)\) sirve como índice de confianza para descartar zonas problemáticas, como reflejos especulares intensos causados por el acero pulido o lecturas con ruido excesivo. A su vez, la radiación ambiental parásita \(A(x, y)\) se filtra promediando las tomas:

\begin{equation}
	A(x, y) = \frac{A_0(x, y) + A_1(x, y) + A_2(x, y) + A_3(x, y)}{4}
\end{equation}

Conocido el ángulo \(\varphi(x, y)\) y con la velocidad de la luz en el aire (\(c \approx \qty{3e8}{\metre\per\second}\)), la distancia real \(R(x, y)\) hasta la superficie se obtiene mediante:

\begin{equation}
	R(x, y) = \frac{c}{4\pi f_{\text{mod}}} \varphi(x, y) = \frac{c \cdot \tau(x, y)}{2}
\end{equation}

El rango de medición queda acotado por \(R_{\text{max}} = \frac{c}{2 f_{\text{mod}}}\).
Finalmente, a partir de la distancia \(R(x, y)\) y los parámetros de calibración del lente, los datos se proyectan a un espacio cartesiano como una nube de puntos 3D \(\mathbf{P}(x, y) = [X, Y, Z]^T\), permitiendo mapear pandeos y variaciones de planitud de la lámina en tiempo real.



% \section{Estado del arte}
%
% \section{Análisis comparativo}
